% Generated draft: requires mathematical and source review.
\documentclass[11pt]{article}
\usepackage[margin=25mm]{geometry}
\usepackage{fontspec}
\setmainfont{Latin Modern Roman}
\usepackage{amsmath,amssymb,mathtools,unicode-math}
\setmathfont{Latin Modern Math}
\usepackage{enumitem,needspace,xurl,hyperref}
\setlength{\parindent}{0pt}
\setlength{\parskip}{5pt}
\setlength{\emergencystretch}{3em}
\newcommand{\N}{\mathbb{N}}
\newcommand{\Z}{\mathbb{Z}}
\newcommand{\Q}{\mathbb{Q}}
\newcommand{\R}{\mathbb{R}}
\newcommand{\C}{\mathbb{C}}
\newcommand{\F}{\mathbb{F}}
\newcommand{\A}{\mathbb{A}}
\newcommand{\PP}{\mathbb{P}}
\newcommand{\E}{\mathbb{E}}
\newcommand{\eps}{\varepsilon}
\newcommand{\dd}{\,\mathrm{d}}
\begin{document}
\section*{UnsolvedMath ID: 20002367}
\textbf{Source problem code:} catalogue entry
\subsection{Definitions and mathematical statement}
Let $K$ be a number field, $\mathcal O_K$ its ring of integers, $d=[K:\Q]$, and $\mathcal O_K^\times$ its unit group. Use $N(a)=|N_{K/\Q}(a)|$ for the absolute field norm; the original writes $N(a)$ without specifying the sign convention. A subset $A\subseteq\mathcal O_K$ is division-ample if it is Diophantine over $\mathcal O_K$, every $x\in\mathcal O_K$ divides some $a\in A$, and
\[
\exists\ell\in\Z_{>0}\ \forall a\in A\ \exists\widetilde a\in\Z:
\quad \widetilde a\mid a\text{ in }\mathcal O_K,\qquad N(a)\leq|\widetilde a|^\ell.
\]
Here Diophantine means existentially definable by finitely many polynomial equations over $\mathcal O_K$, with fixed coefficients permitted. The source notes that for $A\subseteq\Z$ the norm condition follows by taking $\widetilde a=a$ and $\ell=d$.

Question 31 (Cornelissen) asks for such an $A$ satisfying
\[
\forall S\subseteq\mathcal O_K\text{ finite},\qquad
A\not\subseteq\mathcal O_K^\times\cdot(\Z\cup S),
\]
where the product consists of all $u z$ with $u$ a unit and $z\in\Z\cup S$. The source attributes to Cornelissen--Pheidas--Zahidi the implication that existence of a division-ample set, together with an elliptic curve of rank one over $K$, yields a negative answer to Hilbert's tenth problem over $\mathcal O_K$. That is a conditional consequence, not an assumption that every $K$ has such an elliptic curve.

The aid reports the following classification \emph{assuming that $\Z$ is Diophantine over $\mathcal O_K$}. If $d>1$, choose $x_0\in\mathcal O_K\setminus\mathcal O_K^\times\Z$, effectively enumerate $\mathcal O_K=(b_j)_{j\geq0}$, and put
\[
n_j=\max\{1,N(b_j)\}.
\]
Choose distinct increasing positive integers $m_j\geq\max\{1,N(x_0)\}$, and set
\[
A=\{n_jb_j:j\geq0\}\ \cup\ \{m_jx_0:j\geq0\}.
\]
The aid asserts this is computably enumerable and therefore Diophantine under the stated $\Z$-definability hypothesis, using MRDP in integral coordinates. Divisibility density follows from the first family. For its elements take $\widetilde a=n_j$; for the second family take $\widetilde a=m_j$. Both satisfy the norm bound with $\ell=d+1$. The increasing norms of $m_jx_0$ give infinitely many distinct unit orbits, none meeting $\Z$, so no finite $S$ covers $A$ as above. An enumeration including $b_j=0$ includes $0\in A$, accounting also for the literal divisibility demand at $x=0$. When $K=\Q$, the required nonconcentration is impossible because $\mathcal O_K=\Z$ already.

The summary calls the $\Z$-definability hypothesis a now-proved theorem and describes this as a full post-2025 classification, but supplies no author, theorem statement, or authoritative citation for that decisive input. Accordingly the construction and classification are stated conditionally here; the supplied solved labels are not independent verification of that literature claim.
\subsection{Short English statement}
Cornelissen asks for a division-ample set that cannot be covered by integer unit orbits and finitely many additional unit orbits. The supplied two-family construction gives such a set in every proper number field, conditional on a Diophantine definition of the integers in its ring of integers. The rational field cannot satisfy the nonconcentration requirement; the claimed universal definability input needs an authoritative reference.
\subsection{Sources}
\begin{itemize}
\item Local UnsolvedMath registry, permanent ID $20002367$, \texttt{catalogue entry}; \texttt{aim-logic-notes.json}, \texttt{notes[142]}.
\item American Institute of Mathematics, \emph{Extensions of Hilbert's tenth problem}, Question 31 (Cornelissen), \url{https://aimath.org/WWN/hilberts10th/hilberts10th.pdf}. Supplied verification reports page-6 content-stream inspection and restoration of $\widetilde a$, $\Z$, and the unit-group superscript. No recovered passage follows the verification's introductory line.
\item The source attributes the division-ample/rank-one implication to Cornelissen--Pheidas--Zahidi, without a full citation. The aid names MRDP and an unspecified post-2025 theorem that $\Z$ is Diophantine over every number-field ring of integers; no authoritative reference for the latter is supplied.
\item Research attempt 1: full solution, problem solved at run, rigor label ``proved,'' novelty confidence low. Registry classification solved in literature, status solved, statement status exact are supplied labels, qualified by the missing decisive citation.
\end{itemize}
\end{document}
